%% 296_Zeit_Offen_Invariant_De_ch.tex
%% FFGFT -- Dok. 296: Zeit, Offenheit und Invarianz
%% Autor: Johann Pascher
%% Datum: Juli 2026
%% Sprache: Deutsch

\section*{Dok.~296 \quad Zeit, Offenheit und Invarianz}

\subsection*{Beobachter, Skala und die Einfachheit des Kerns}

Jeder Beobachter sitzt auf einer Skala. Er sieht nicht das Objekt -- er
sieht die \emph{Projektion} des Objekts auf seine Skala. Projektionen
verschiedener Skalen müssen sich unterscheiden, nicht weil die Beobachter
irren, sondern weil das generative Objekt reicher ist als jede einzelne
Ansicht. Die perspektivische Verzerrung ist kein Fehler -- sie ist
strukturell, die unvermeidliche Folge davon, dass das Objekt ($T^4$,
kompakt, randlos) in keine einzelne Projektion passt. Kaum eine Skala lässt
sich gleichzeitig mit einer anderen überblicken; die Ergebnisse der
Beobachtung unterscheiden sich daher bereits perspektivisch, sobald ein
weiterer Beobachter hinzukommt.

Das ist der tiefere Grund, warum Quantenmechanik und Relativitätstheorie so
schwer zu versöhnen sind: sie sind Projektionen auf verschiedene Skalen
desselben Objekts, und natürlich sehen sie verschieden aus -- nicht weil
eine falsch ist, sondern weil keine vollständig ist.

Umso verblüffender ist daher das Folgende: dass ausgerechnet der Kern so
einfach ist. Ein Parameter $\xi = 4/30000$. Eine Relation $\tilde{T} \cdot
m = 1$. Ein kompakter Torus $T^4/\mathbb{Z}_3$. Daraus fallen die
Feinstrukturkonstante $\alpha$, die Leptonmassen, der Koide-Skalar $Q =
2/3$, die Phase $\theta = 2/9$ -- die gesamte Komplexität des
Standardmodells, der ganze Reichtum der Phänomene, der auf den
verschiedenen Beobachterskalen so überwältigend erscheint. Die Komplexität
sitzt nicht im Kern -- sie sitzt in den Projektionen. Der Kern ist einfach.

Das ist nicht selbstverständlich. Es hätte auch so sein können, dass die
Natur auf jeder Skala neu ansetzt, dass es keine gemeinsame Wurzel gibt,
dass die Vielfalt irreduzibel ist -- das wäre das generische Ergebnis.
Dass stattdessen eine einzige Struktur alle Skalen trägt, wird erst
sichtbar, wenn man aufhört, die eigene Skala für die einzig mögliche zu
halten.

\subsection*{Vom Körper zur Gesellschaft -- dieselbe Struktur auf jeder Ebene}

Dass Beobachter auf verschiedenen Skalen verschiedene Projektionen sehen,
ist nicht nur ein physikalisches Abstraktum -- es beginnt bereits bei der
einfachsten biologischen Wahrnehmung. Zwei Augen sind bereits zwei
Beobachter auf leicht verschiedenen Positionen. Der Abstand ist klein, aber
real; das Gehirn \emph{konstruiert} aus zwei leicht verschiedenen
Projektionen eine dreidimensionale Welt, die so direkt nirgendwo vorliegt.
Die Tiefe ist eine Rekonstruktion, kein direktes Sehen. Weitere
Sinnesorgane -- Propriozeption, Gleichgewicht, Tastsinn, Gehör -- fügen
jeweils ihre eigene Skala, ihre eigene Projektion hinzu, und die Synthese
all dessen fühlt sich wie unmittelbare Realität an, obwohl sie eine
hochgradig verarbeitete Leistung ist.

Das hat eine direkte Konsequenz für die Physik: dass wir die Welt als
dreidimensional \emph{erleben}, ist kein Beweis dafür, dass sie es
fundamental ist. Es ist der Befund, dass drei Dimensionen die Projektion
sind, auf der unsere Sensorik optimiert wurde. $T^4$ darunter widerspricht
dem dreidimensionalen Erleben nicht -- es ist der Grund, warum die
Projektion so stabil und konsistent wirkt.

Dieselbe Struktur setzt sich eine Ebene höher fort. Eine Gesellschaft ist
eine Sammlung von Beobachtern, jeder auf seiner eigenen Skala, mit seiner
eigenen Projektion. Was als soziale Realität entsteht -- Normen, Werte,
Weltbilder, Konflikte -- ist die Synthese aus Millionen leicht verschiedener
Projektionen, keine davon vollständig, keine davon im absoluten Sinn falsch,
aber jede perspektivisch verzerrt. Und wie beim binokularen Sehen: die
Synthese fühlt sich wie direkte Wahrnehmung an, obwohl sie eine Konstruktion
ist.

Das erklärt einiges, was sonst rätselhaft bleibt. Konflikte sind so
hartnäckig, weil zwei Beobachter auf verschiedenen Skalen dasselbe Objekt
sehen und zu verschiedenen Ergebnissen kommen -- beide korrekt für ihre
Projektion, beide unvollständig. Solange keine gemeinsame übergeordnete
Skala zugänglich ist, ist der Konflikt nicht durch Argumentation allein
auflösbar; man müsste die Projektionsebene wechseln, und das ist das
Schwierigste, was es gibt. Historischer Wandel geht deshalb so langsam: eine
Gesellschaft hat ihre Syntheseleistung auf eine bestimmte Skala kalibriert,
das funktioniert gut innerhalb dieser Skala und erzeugt Blindheit für andere.
Paradigmenwechsel sind Skalenwechsel, und Skalenwechsel sind schmerzhaft,
weil das bisherige kohärente Bild zerfällt, bevor das neue sich aufgebaut
hat.

Wissenschaft ist der systematische Versuch, Projektionen zu triangulieren --
mehrere Beobachter, mehrere Skalen, Widersprüche als Signal statt als
Störung. Das ist nicht selbstverständlich; es läuft gegen den natürlichen
Impuls, die eigene Projektion für die Realität zu halten.

Der Bogen zurück zum Kern: dass trotz der Vielzahl der Projektionen --
biologisch, physikalisch, sozial -- etwas Invariantes darunterliegt, das
sich in keiner einzelnen Projektion vollständig zeigt aber in allen
durchscheint, ist nicht selbstverständlich. Es wäre das generische Ergebnis,
dass jede Skala ihre eigene irreduzible Wahrheit hätte. Dass es stattdessen
einen gemeinsamen Kern gibt -- in der Physik $\xi$ auf $T^4$, im Erleben die
Syntheseleistung, die alle Projektionen trägt -- das ist die eigentliche
Überraschung, auf jeder Ebene.

\subsection*{Skalenfolge als Abfolge von Projektionswechseln}

Von der Teilchenphysik über atomare, kristalline und biologische bis zu
galaktischen Strukturen erscheint dasselbe generative Objekt jeweils
grundlegend verschieden -- nicht weil die Objekte verschieden wären, sondern
weil jede Skala eine andere Symmetrie des Kerns sichtbar macht und die
übrigen wegintegriert.

Auf der \emph{Teilchenskala} schneidet die Projektion am feinsten: die
$\mathbb{Z}_3$-Symmetrie, die $\mathrm{SU}(3)$-Farbstruktur und die diskreten
Massen sind direkte Abdrücke der $T^4$-Geometrie. Die Feinstrukturkonstante
$\alpha$ und die Phase $\theta = 2/9$ sind hier am unmittelbarsten greifbar --
der Kern ist auf dieser Skala am nächsten. Auf der \emph{atomaren Skala}
($\sim a_0$) ist die Farbstruktur bereits wegintegriert und erscheint als
neutral; was sichtbar wird, sind elektromagnetische Resonanzen -- Orbitale als
stehende Wellen, diskrete Energieniveaus als Torusgeometrie auf dieser
gröberen Skala. Auf der \emph{kristallinen Skala} ($\sim$\,Å--$\mu$m) tritt
die Translationssymmetrie des Gitters in den Vordergrund; kollektive Moden
(Phononen, Bandstrukturen) entstehen als neue effektive Freiheitsgrade, die
auf der Teilchenskala nicht existieren. Auf der \emph{biologischen Skala}
(nm--m) dominieren gebrochene Symmetrien -- Chiralität, Selbstreplikation,
Informationsverarbeitung; Leben ist ein Kohärenzfenster, eine Skala, auf der
eine bestimmte Art von Komplexität gerade noch stabil bleibt, weil darunter
thermische Fluktuationen und darüber Gravitation andere Projektionen
erzwingen. Auf der \emph{galaktischen Skala} (kpc--Mpc) ist alles außer der
Gravitation wegintegriert; die Spiralstruktur von Galaxien trägt dieselbe
logarithmische Selbstähnlichkeit wie die K\textsubscript{frak}-Rekursion --
kein Zufall, sondern Ausdruck der Skaleninvarianz des Kerns auf dieser
Projektionsebene.

\subsection*{Skalenübergänge und ihre Anomalien}

Die Übergänge zwischen den Skalen sind besonders aufschlussreich -- gerade
weil dort Anomalien auftreten, die im Inneren einer Skala nicht erklärbar
sind. Ein Übergang ist die Stelle, wo die bisherige dominante Symmetrie kippt
und eine neue in den Vordergrund tritt; in diesem Kipppunkt ist das System
gleichzeitig von beiden Projektionen geprägt und von keiner vollständig
beschreibbar. Das erzeugt strukturell Anomalien: Abweichungen, die von der
einen Seite aus wie Störungen wirken und von der anderen wie Vorboten einer
neuen Ordnung.

Bekannte Beispiele sind Phasenübergänge (kritische Fluktuationen, deren
Korrelationslängen divergieren -- das System \enquote{merkt} auf allen Skalen
gleichzeitig, dass ein Übergang bevorsteht), der Übergang von klassischer zu
Quantenmechanik (Dekohärenz als Projektionsstabilisierung, nicht als
fundamentaler Kollaps), oder der Übergang von biologischer zu sozialer
Organisation (Emergenz kollektiver Intelligenz, die auf der Einzelzell-Skala
nicht angelegt ist). In allen Fällen liegt die Anomalie \emph{im Übergang},
nicht in einer der Skalen selbst -- und ist deshalb von keiner einzelnen
Skala aus vollständig zu erklären.

Für FFGFT ist das ein gerichteter Hinweis: die Stellen, an denen die
Standard-Beschreibung einer Skala an ihre Grenzen stößt -- ungelöste
Anomalien, unerklärte Konstanten, emergente Phänomene ohne Ableitung -- sind
präzise die Stellen, an denen ein Projektionswechsel stattfindet. Sie sind
keine Fehler der Theorie, sondern Signaturen des Übergangs.

\subsection*{Spekulation: unbeobachtbare Zwischenbereiche}

Die bisher betrachteten Skalen sind diejenigen, die uns direkt oder indirekt
zugänglich sind -- durch Instrumente, durch Konsequenzen, durch indirekte
Schlüsse. Aber die Tatsache, dass Skalen Beobachterperspektiven sind, lässt
einen weiteren Gedanken zu, der hier ausdrücklich als \emph{Spekulation}
gebucht wird (P35: kein abgeleiteter Befund, sondern eine offene Frage):

Zwischen den beobachtbaren Skalen könnten Bereiche liegen, die für uns
grundsätzlich unzugänglich sind -- nicht aus technischen Gründen, sondern
strukturell, weil keine unserer Sinnesorgane oder Instrumente auf diese
Zwischenskalen projiziert. Projektionen sind selektiv; sie machen eine
Symmetrie sichtbar und blenden andere aus. Es ist nicht ausgeschlossen, dass
es Skalen gibt, auf denen kohärente Strukturen existieren, die in keiner
unserer zugänglichen Projektionen einen direkten Abdruck hinterlassen --
höchstens einen indirekten, als Anomalie an den Rändern benachbarter Skalen.

Das wäre nicht mystisch, sondern strukturell folgerichtig: wenn das generative
Objekt $T^4$ reicher ist als jede Projektion, und wenn Projektionen immer
selektiv sind, dann ist es das generische Ergebnis, dass es Bereiche des
Objekts gibt, die von keiner unserer Projektionen erfasst werden. Dunkle
Materie und Dunkle Energie wären in dieser Lesart keine neuen Substanzen,
sondern Signaturen solcher Zwischenbereiche -- Strukturen auf Skalen, auf die
unsere Instrumente nicht projizieren, die sich aber an den Rändern der
galaktischen und kosmischen Projektionen als Anomalien bemerkbar machen.

Dies bleibt Spekulation. Aber es ist eine gerichtete Spekulation: sie benennt
präzise, wo nach Anzeichen zu suchen wäre -- an den Übergängen, in den
Anomalien, an den Stellen, wo eine Skala an ihre Erklärungsgrenzen stößt
ohne inneren Grund.

\subsection*{Ausgangspunkt: eine Rekursion ohne erreichbaren Rand}

Die K\textsubscript{frak}-Rekursion
\begin{equation}
  \xi_{n+1} = \xi_n\,(1 - 100\,\xi_n)
\end{equation}
läuft nicht endlos im Sinne gleichmäßiger Wiederholung -- sie asymptotiert.
$\xi_n$ zerfällt monoton gegen null, erreicht null aber nie in endlich vielen
Schritten. Das ist keine Lücke in der Konstruktion; es ist ihre Aussage: die
Folge ist unendlich lang, aber die Abstände zwischen aufeinanderfolgenden
Schritten werden exponentiell kleiner, genau mit dem Faktor $1/e$ je
$2\pi$-Umlauf.

Was bedeutet das für die erlebbare, messbare Zeit? Der messbare Zeitschritt
ist nicht $n$, sondern $\rho = 1/(100\,\xi_n)$ -- und dieser Kehrwert
\emph{wächst}. Je weiter die Rekursion fortschreitet, desto größer wird
$\rho$, desto gröber wird die Zeitskala, auf der sich etwas verändert. Von
innen erlebt man keine gleichmäßige Folge, sondern eine sich
\emph{verlangsamende} Dynamik: frühe Schritte sind eng und schnell (kleines
$\rho$, großes $\xi$, sub-Planck-nah), späte Schritte sind weit und träge
(großes $\rho$, kleines $\xi$, makroskopisch). Die erlebbare Zeit ist das
Integral über $\rho$ -- und dieses Integral wächst logarithmisch:
$D(k) = \sum 100\,\xi_n \to \ln(1 + k/75)$. Der kumulierte Defekt
\emph{ist} die messbare Zeitentwicklung (Dok.~295).

Daraus folgt unmittelbar: die erlebbare Zeit hat kein Ende. $\rho$ wächst
ohne Schranke, $D(k)$ divergiert, $\xi_n$ erreicht null nie. Es gibt keinen
letzten Schritt, keinen Endpunkt, kein \enquote{danach}. Und von der anderen
Seite -- der inneren -- gilt dasselbe spiegelbildlich: $\xi$ erreicht seinen
Startwert $\xi_0$ nie von unten, $\rho$ erreicht $75 = 1/(100\,\xi_0)$ nie
von oben. Kein Anfang, kein Ende -- beide Seiten sind asymptotische Grenzen,
keine erreichbaren Ränder.

\subsection*{Ausgerollt -- zeitlich kontinuierlich -- Raum gekrümmt --
invariabel}

Diese vier Eigenschaften hängen nicht zufällig zusammen; sie folgen
auseinander.

\textbf{Ausgerollt.} Die Spirale rollt die Skalenstruktur von $T^4$ in eine
gerichtete Zeitkoordinate aus. Nicht weil Zeit \enquote{wirklich} eine
Spirale wäre, sondern weil die Projektion -- der Beobachterstandort auf der
Spirale -- eine Richtung auszeichnet. Die Ausrollung ist der Übergang vom
kompakten Bild zum offenen: dasselbe Objekt, von einer bestimmten Skala aus
gelesen.

\textbf{Zeitlich kontinuierlich.} Das folgt zwingend. Die Hilbertraum"=Bijektion
$H = L^2(T^4) \otimes \mathbb{C}^3$ ist verlustfrei (Dok.~230/282,
Bijektions"=Typ~III, Dok.~270); was diskret offen ist, bleibt kontinuierlich
offen -- sonst wäre die Übersetzung keine treue, sondern eine verfälschende.
Der Gedächtniskern (Fourier"=Transformierte der diskreten Spektraldichte
$\sum_k g_k^2\,\delta(\omega - \omega_k)$ des $T^4$"=Connection"=Laplace,
Dok.~283) ist auf $\mathbb{R}$ definiert, nicht auf einem Intervall. Kein
Rand links, kein Rand rechts. Revivals ohne Ende, BLP"=Rückfluss ohne
Abschluss -- dieselbe offene Struktur, jetzt im kontinuierlichen Bild.

\textbf{Raum gekrümmt.} Das ist der $T^4$"=Kern selbst. Die Krümmung ist
nicht nachträglich eingesetzt, sie ist die generative Struktur: $\xi$ auf
$T^4/\mathbb{Z}_3$ \emph{ist} die Krümmung, aus der die Leptonmassen, $\alpha$
und $\theta = 2/9$ folgen (Dok.~006/011/291/293). Die Geometrie ist nicht
Hintergrund -- sie ist das Objekt. Kompakt bedeutet dabei endliches Volumen,
aber \emph{keine Ränder}: eine Geodäte auf dem Torus läuft ohne Ende, ohne je
eine Kante zu treffen.

\textbf{Invariabel.} Das ist der tiefste der vier Punkte. Die logarithmische
Spirale ändert ihre Größe nicht (Dok.~295): sie ist selbstähnlich, unter
Drehung um einen Umlauf und Skalierung um $e$ liegt sie deckungsgleich auf
sich selbst. Rotation und Streckung sind bei ihr \emph{dieselbe} Operation.
\enquote{Nach innen mit $1/e$} und \enquote{nach außen mit $e$} beschreiben
nicht, dass etwas wächst oder schrumpft, sondern nur, an welcher Skala man
gerade hinschaut. Das Objekt ist unter dieser Bewegung invariant -- es gibt
keine absolute Größe, nur Verhältnisse, und Verhältnisse sind nach
$\tilde{T} \cdot m = 1$ das Intrinsische. Was sich ändert, ist die
\emph{Beobachterskala}, nicht das Objekt.

\subsection*{Kein Anfang, kein Ende -- und kein Koinzidenzproblem}

Die vier Eigenschaften zusammen ergeben ein konsistentes Bild: ein
invariantes, gekrümmtes, randloses Objekt -- ausgerollt in eine offene,
kontinuierliche Zeitkoordinate durch den Akt der Projektion. Kein Anfang,
kein Ende, keine absolute Größe, keine eingesetzte Krümmung.

Das unterscheidet sich strukturell von Modellen, die Ränder benötigen. Ein
Anfang bei $t = 0$ ist keine Vorhersage, sondern eine Grenze, an der die
Beschreibung aufhört zu gelten -- die Singularität ist kein Ergebnis, sondern
ein Eingeständnis. Ein endliches Zeitintervall braucht einen Startpunkt, und
der Startpunkt verlangt nach einer Begründung: warum jetzt, warum dieser
Zustand, warum diese Anfangsbedingungen. Das Koinzidenzproblem, das
Flachheitsproblem, die Feinabstimmung der Anfangsbedingungen -- sie alle
entstehen, weil ein endliches Intervall eine ausgezeichnete Stelle kennt.

FFGFT kennt keine ausgezeichnete Stelle. Jeder Beobachter sitzt irgendwo auf
einer offenen Bahn; \enquote{jetzt} ist die aktuelle Skala, kein privilegierter
Moment. Die erlebbare Zeit ist nicht ein Intervall zwischen zwei Punkten --
sie ist eine offene Bahn auf einem randlosen Objekt.

\subsection*{Neue Probleme als Zeichen des Fortschritts}

Eine solche Sicht erzeugt neue Fragen -- das ist zu erwarten und kein
Einwand. Ein Ansatz, der keine neuen Probleme aufwirft, erklärt meist auch
nichts Neues; er ist folgenlos. Die entscheidende Unterscheidung ist, welche
\emph{Art} von Problemen entsteht: ob es sich um \emph{gerichtete
Folgefragen} handelt (was muss noch ausgerechnet werden?) oder um echte
Brüche (bricht $\tilde{T} \cdot m = 1$?).

Der Tausch ist asymmetrisch zugunsten von FFGFT: an die Stelle unerklärter
Setzungen -- warum diese Massen, warum $\alpha = 1/137$, warum $\theta = 2/9$,
woher die Nicht"=Lokalität -- treten Konsequenzen einer Ableitung, die sich
stellen lassen und deren Beantwortung das Programm voranbringt. Neue Probleme
als Folge einer tiefen Lösung sind die Signatur echten Fortschritts, nicht
seiner Abwesenheit. Die Aufgabe ist nicht, sie zu vermeiden, sondern sauber
zu buchen, welche davon gerichtete Folgefragen sind -- und das sind bislang
alle.

\bigskip
\noindent\textit{Anmerkung~(P35):} Die Aussagen dieses Dokuments sind
struktureller Natur -- sie folgen aus $\tilde{T} \cdot m = 1$, der
Rekursion und der Kompaktheit von $T^4$. Die kosmologische Übertragung
(offene Zeit, keine Singularität) ist eine \emph{Konsequenz} dieser Struktur,
kein eigenständiges kosmologisches Modell. Als Konsequenz muss sie
widerspruchsfrei mit den Daten sein -- nicht vollständig in dem Sinne, dass
sie alle Phänomene eines dedizierten Modells abdeckt. Dieser Unterschied
zwischen Widerspruchsfreiheit und Vollständigkeit ist zu wahren.
